\section{Conclusion}
\label{sec:conclusion}

Under physical-device heterogeneity in N-BaIoT, threshold scope alone strongly changes
the distribution of false alarms across clients: per-client $p_{95}$ calibration (B2)
reduced CV(FPR) from $1.017$ to $0.299$ under the fixed FedAvg AE.
Cluster-mean calibration (B4) recovers $\approx$52\% of B2's improvement without
requiring a device taxonomy, at the cost of transmitting distributional metadata that
is not private; family-mean calibration (B3) performs poorly because coarse device-type
labels do not align with reconstruction-error calibration structure.
Under near-homogeneous conditions (Regime~B file-level CICIoT2023, Dirichlet near-IID)
no dispersion reduction is observed---an applicability boundary, not a general
recommendation against personalization.
B2 is therefore an FPR-equity intervention, not a universal detection-quality improvement:
lower-tail P10 Macro-F1 degradation concentrates in the low-separability Ennio Doorbell
client. Practitioner guidance is therefore two-step: first measure benign calibration-error heterogeneity, using pairwise distribution divergence or dispersion of local calibration thresholds; then choose the threshold policy according to that heterogeneity and deployment constraints. B1 is appropriate when clients are near-homogeneous or per-client state is undesirable; B2 should be considered when local calibration is available and a shared-threshold pilot shows high relative FPR dispersion, for example CV(FPR) near or above~1 with full calibration coverage; B4 is a taxonomy-free compromise when grouping is acceptable, while accounting for metadata leakage from shared distributional fingerprints.

\textbf{Limitations and future work.}
Regime~A uses only 9 physical devices (B4 exploratory at this scale);
Regime~B uses randomized CICIoT2023 file-level pseudo-clients;
$E{=}1$ with full participation limits local drift; five seeds limit statistical power;
no adversarial robustness, formal privacy, or hardware evaluation is provided.
Future work covers larger fleets, $E>1$ sensitivity, privacy-preserving grouping,
model-level personalization, adversarial robustness, and deployment validation under
concept drift.

\textbf{Availability.} Artifact: Zenodo \nolinkurl{10.5281/zenodo.20315654}; release v1.0.0; license and reproduction commands are in the repository.
